\documentclass[brazil,12pt,a4paper,openany]{book}
\usepackage{anysize}
%\usepackage[latin1]{inputenc}
\usepackage[utf8]{inputenc}
\usepackage[portuguese,brazilian]{babel}
\usepackage[portuguese]{babel}
\usepackage[T1]{fontenc}
\usepackage[ruled,chapter]{algorithm}
\usepackage{makeidx}
\usepackage{subfig}
\usepackage{amsmath}
\usepackage{amssymb}
\usepackage{amsthm}
\usepackage{mathrsfs}
\usepackage{graphicx}
\usepackage{fancyhdr}
\usepackage{array}
\usepackage{geometry}
\usepackage{simplewick}
\usepackage{latexsym}
\usepackage[all]{xy}
\usepackage{euscript}
\usepackage{enumerate}
\usepackage{dsfont}
\usepackage{slashed}
\usepackage[plainpages=false, pagebackref=true, pdftex]{hyperref}
\usepackage{braket}

\usepackage{textcomp}
\usepackage{siunitx}
\usepackage{gensymb}

\newcommand{\G}{\mathcal G}

\usepackage{pdfpages}
\usepackage{indentfirst}
\usepackage[intoc]{nomencl}

\makenomenclature
\renewcommand{\nomname}{Lista de Símbolos}

\makeindex

\marginsize{25mm}{25mm}{25mm}{25mm}

\hypersetup{colorlinks=true, breaklinks=true, linkcolor=black, citecolor=black, urlcolor=black}

\begin{document}

\setlength{\abovecaptionskip}{0.0cm}
\setlength{\belowcaptionskip}{0.0cm}
\setlength{\baselineskip}{24pt}

\pagestyle{fancy}
\lhead{}
\chead{}
\rhead{\thepage}
\lfoot{}
\cfoot{}
\rfoot{}

\fancypagestyle{plain}
{
	\fancyhf{}
	\lhead{}
	\chead{}
	\rhead{\thepage}
	\lfoot{}
	\cfoot{}
	\rfoot{}
}

\renewcommand{\headrulewidth}{0pt}



%%%%%%%%%%%%%%%%%%%%%%%%%%%%%%%%% Capa %%%%%%%%%%%%%%%%%%%%%%%%%%%%%%%%%%

\frontmatter 

\thispagestyle{empty}

\begin{figure}[h]
	\includegraphics[scale=0.8]{fig_doc/logo_ufrj.jpg}
\end{figure}

\vspace{15pt}

\begin{center}

\textbf{UNIVERSIDADE FEDERAL DO RIO DE JANEIRO}

\textbf{INSTITUTO DE FÍSICA}

\vspace{30pt}

{\Large \bf Título da tese ou dissertação}

\vspace{25pt}

{\large \bf Nome completo do candidato}

\vspace{35pt}

\begin{flushright}
\parbox{10.3cm}{Dissertação de Mestrado / Tese de Doutorado apresentada ao Programa de Pós-Graduação em Física do Instituto de Física da Universidade Federal do Rio de Janeiro - UFRJ, como parte dos requisitos necessários à obtenção do título de Mestre / Doutor em Ciências (Física).

\vspace{18pt}

{\large \bf Orientador: Nome completo do orientador}

\vspace{12pt}

{\large \bf Coorientador: Nome completo do coorientador}}
\end{flushright}

\vspace{90pt}

\textbf{Rio de Janeiro}

\textbf{\today}

\end{center}

%%%%%%%%%%%%%%%%%%%%%%%%%%%%%%%%%%%%%%%%%%%%%%%%%%%%%%%%%%%%%%%%%%%%%%%%%



%%%%%%%%%%%%%%%%%%%% Folha de Assinaturas da Banca %%%%%%%%%%%%%%%%%%%%%%

\newpage

\thispagestyle{empty}

\noindent
\setlength{\headheight}{14.49998pt}

%\begin{figure}[ht]
%  \begin{center}
%	\includegraphics[scale=0.742]{fig_doc/Folha_Ata.pdf}
%  \end{center}
%\end{figure}

\clearpage

%%%%%%%%%%%%%%%%%%%%%%%%%%%%%%%%%%%%%%%%%%%%%%%%%%%%%%%%%%%%%%%%%%%%%%%%%



%%%%%%%%%%%%%%%%%%%%%%%% Ficha Catalográfica %%%%%%%%%%%%%%%%%%%%%%%%%%%%

\newpage
\mbox{}\vspace{5cm}

\begin{center}
\begin{tabular}{|c|}
\hline\\
\begin{minipage}{15cm}
\rule{15cm}{0pt} \flushright


\begin{minipage}{12.5cm}

[Somente último sobrenome], [Primeiro nome + sobrenome do meio]

% VER http://biblioteca.if.ufrj.br/wp-uploads/2014/09/tutorial-ficha-catalografica.pdf
% VER TAMBÉM http://biblioteca.if.ufrj.br/servicos-produtos/ficha-catalografica/
\hspace{-2cm} CODIGO CUTTER 
\hspace{1,59cm} Título da tese / dissertação / Nome do candidato -- Rio de Janeiro, Ano.

\hspace{1cm} Número de folhas

\hspace{2cm} 

\hspace{1cm} Orientador: .

\hspace{1cm} Coorientador: .

\hspace{1cm} Tese/Dissertação (doutorado/mestrado) - Universidade Federal do \\Rio de Janeiro, Instituto de Física, Programa de \\Pós-Graduação em Física, ano.

\hspace{1cm} Referências Bibliográficas: f. [primeira folha das referências biblio]-[última].

\hspace{1cm} 1. palavra-chave 1. 2. palavra-chave 2.\\ 3. palavra-chave 3. 4. palavra-chave 4. 5. palavra-chave 5. \\6. palavra-chave 6. I. [último sobrenome do orientador], [nome + sobrenome do meio do orientador], \\orient. II. [último sobrenome do coorientador], [nome + sobrenome do meio do coorientador], coorient. III. Universidade \\Federal do Rio de Janeiro, Instituto de Física, Programa \\de Pós-graduação em Física. IV. Título

\bigskip
\end{minipage}
\end{minipage}
\\
\hline
\end{tabular}
\end{center}

%%%%%%%%%%%%%%%%%%%%%%%%%%%%%%%%%%%%%%%%%%%%%%%%%%%%%%%%%%%%%%%%%%%%%%%%%



%%%%%%%%%%%%%%%%%%%%%%%%%%%%%%%% Resumo %%%%%%%%%%%%%%%%%%%%%%%%%%%%%%%%%

\newpage

\noindent

\vspace*{20pt}
\begin{center}
{\LARGE\bf Resumo}\\
\vspace{15pt}
{\Large\bf Título da Tese/Dissertação}\\
\vspace{6pt}
{\bf Nome do candidato}\\
\vspace{12pt}
{\bf Orientador: }\\
{\bf Coorientador: }\\
\vspace{20pt}
\parbox{14cm}{Resumo da Dissertação de Mestrado apresentada ao Programa de Pós-Graduação em Física do Instituto de Física da Universidade Federal do Rio de Janeiro - UFRJ, como parte dos requisitos necessários à obtenção do título de Mestre em Ciências (Física).}
\end{center}
\vspace*{35pt}

TEXTO DO RESUMO

\vspace{15pt}

\textbf{Palavras-chave:} Palavra-chave 1,  Palavra-chave 2, etc.

%%%%%%%%%%%%%%%%%%%%%%%%%%%%%%%%%%%%%%%%%%%%%%%%%%%%%%%%%%%%%%%%%%%%%%%%%



%%%%%%%%%%%%%%%%%%%%%%%%%%%%%%% Abstract %%%%%%%%%%%%%%%%%%%%%%%%%%%%%%%%

\newpage

\noindent

\vspace*{20pt}
\begin{center}
{\LARGE\bf Abstract}\\
\vspace{15pt}
{\Large\bf Título em inglês}\\
\vspace{6pt}
{\bf Nome do candidato}\\
\vspace{12pt}
{\bf Orientador: }\\
{\bf Coorientador: }\\
\vspace{20pt}
\parbox{14cm}{\emph{Abstract} da Dissertação de Mestrado apresentada ao Programa de Pós-Graduação em Física do Instituto de Física da Universidade Federal do Rio de Janeiro - UFRJ, como parte dos requisitos necessários à obtenção do título de Mestre em Ciências (Física).} % ATENÇÃO: esse trecho fica em português mesmo!
\end{center}
\vspace*{35pt}

TEXTO DO RESUMO EM INGLÊS

\vspace{15pt}

\textbf{Keywords:} Palavras-chave em inglês.

%%%%%%%%%%%%%%%%%%%%%%%%%%%%%%%%%%%%%%%%%%%%%%%%%%%%%%%%%%%%%%%%%%%%%%%%%



%%%%%%%%%%%%%%%%%%%%%%%%%%%% Agradecimentos %%%%%%%%%%%%%%%%%%%%%%%%%%%%%

\newpage

\noindent

\vspace*{20pt}

\begin{center}

{\LARGE\bf Agradecimentos}

\end{center}

\vspace*{40pt}

Agradecimentos.
Não esquecer de agradecer à agência pela bolsa concedida se for o caso.


%%%%%%%%%%%%%%%%%%%%%%%%%%%%%%%%%%%%%%%%%%%%%%%%%%%%%%%%%%%%%%%%%%%%%%%%%



%%%%%%%%%%%%%%%%%%%%%%%%%%%%%%%%%%%%%%%%%%%%%%%%%%%%%%%%%%%%%%%%%%%%%%%%%

\newpage
\phantomsection
\addcontentsline{toc}{chapter}{Sumário}
\tableofcontents



\printnomenclature

\newpage
\phantomsection
\addcontentsline{toc}{chapter}{Lista de Figuras}
\listoffigures

\mainmatter
\input{capitulo1.tex}
\input{capitulo2.tex}
\input{capitulo3.tex}
\input{capitulo4.tex}
\input{capitulo5.tex}

%%%%%%%%%%%%%%%%%%%%%%%%%%%%%%%%%%%%%%%%%%%%%%%%%%%%%%%%%%%%%%%%%%%%%%%%%



%%%%%%%%%%%%%%%%%%%%%%%%%%% Bibliografia %%%%%%%%%%%%%%%%%%%%%%%%%%%%%%%%

\newpage
\phantomsection


\bibliographystyle{unsrt}
%\bibliographystyle{plain}
\bibliography{referencias}





%%%%%%%%%%%%%%%%%%%%%%%%%%%%%%%%%%%%%%%%%%%%%%%%%%%%%%%%%%%%%%%%%%%%%%%%



%%%%%%%%%%%%%%%%%%%%%%%%%%%%% Apêndices %%%%%%%%%%%%%%%%%%%%%%%%%%%%%%%%%

\appendix
\input{apendicea.tex}
%\input{apendiceb.tex}

%%%%%%%%%%%%%%%%%%%%%%%%%%%%%%%%%%%%%%%%%%%%%%%%%%%%%%%%%%%%%%%%%%%%%%%%%


\end{document}
